\section{Finite Element Approximation and Convergence Analysis}

In this section, we introduce the finite element approximation of the stochastic advection--diffusion problem and derive a convergence estimate.

\subsection{Finite Element Space}

Let $\mathcal{T}_h$ be a regular triangulation of the domain $\Omega$ with mesh size $h$.

We define the finite element space

\[
V_h
=
\left\{
v_h\in H_0^1(\Omega)
:
v_h|_K \in \mathbb P_1(K),
\quad
\forall K\in\mathcal T_h
\right\}.
\]

Throughout this work, we assume that the mesh family is shape-regular and quasi-uniform.

\subsection{Ritz Projection}

Let

\[
R_h : V \rightarrow V_h
\]

denote the Ritz projection defined by

\[
a(R_hu,v_h)
=
a(u,v_h),
\qquad
\forall v_h\in V_h.
\]

The following approximation properties are classical.

\begin{lemma}
Assume that

\[
u\in H^2(\Omega)\cap H_0^1(\Omega).
\]

Then there exists a constant $C>0$ independent of $h$ such that

\[
\|u-R_hu\|_{L^2(\Omega)}
\le
Ch^2
\|u\|_{H^2(\Omega)},
\]

and

\[
\|\nabla(u-R_hu)\|_{L^2(\Omega)}
\le
Ch
\|u\|_{H^2(\Omega)}.
\]

\end{lemma}

\subsection{Semi-Discrete Formulation}

The semi-discrete finite element approximation seeks

\[
u_h(t)\in V_h
\]

such that

\[
(du_h,v_h)
+
a(u_h,v_h)\,dt
=
(f,v_h)\,dt
+
\sigma(v_h)\,dW_t,
\]

for all

\[
v_h\in V_h.
\]

Let

\[
\{\varphi_i\}_{i=1}^{N_h}
\]

be a basis of $V_h$. Writing

\[
u_h
=
\sum_{j=1}^{N_h}
U_j(t)\varphi_j,
\]

we obtain the matrix system

\[
M\,dU
+
A\,U\,dt
=
F\,dt
+
\sigma G\,dW_t,
\]

where

\[
M_{ij}
=
(\varphi_j,\varphi_i),
\]

is the mass matrix and

\[
A_{ij}
=
a(\varphi_j,\varphi_i),
\]

is the stiffness-advection matrix.

\subsection{Fully Discrete Scheme}

Let

\[
t_n=n\Delta t,
\qquad
n=0,\ldots,N.
\]

Using the implicit Euler scheme, we obtain

\[
M U^{n+1}
+
\Delta t\,A U^{n+1}
=
M U^n
+
\Delta t\,F^n
+
\sigma G \Delta W_n,
\]

where

\[
\Delta W_n
=
W_{t_{n+1}}
-
W_{t_n}.
\]

The Wiener increments satisfy

\[
\Delta W_n
\sim
\mathcal N(0,\Delta t).
\]

\subsection{Discrete Stability}

\begin{theorem}
Assume that the hypotheses of Sections 2 and 3 hold.

Then the fully discrete solution satisfies

\[
\max_{0\le n\le N}
\mathbb E
\left[
\|u_h^n\|^2
\right]
+
\Delta t
\sum_{n=0}^{N}
\mathbb E
\left[
\|\nabla u_h^n\|^2
\right]
\le C,
\]

where the constant $C$ is independent of $h$ and $\Delta t$.
\end{theorem}

\begin{proof}[Sketch of proof]

Choosing

\[
v_h=u_h^{n+1}
\]

in the discrete formulation and using

\[
(a-b)a
=
\frac12
\left(
a^2-b^2+(a-b)^2
\right),
\]

one derives a discrete energy inequality.

Taking expectations and applying a discrete Grönwall lemma yields the result.

\end{proof}

\subsection{Error Decomposition}

Let

\[
e^n
=
u(t_n)-u_h^n.
\]

We decompose the error as

\[
e^n
=
\underbrace{u(t_n)-R_hu(t_n)}_{\eta^n}
+
\underbrace{R_hu(t_n)-u_h^n}_{\theta^n}.
\]

The term $\eta^n$ represents the interpolation error, while $\theta^n$ is the discrete error.

Using the approximation properties of the Ritz projection, we obtain

\[
\|\eta^n\|
\le
Ch^2.
\]

\subsection{Convergence Estimate}

We now state the main convergence result.

\begin{theorem}
Assume that the exact solution satisfies

\[
u
\in
L^2
\bigl(
\Omega_P;
H^2(\Omega)
\bigr),
\]

and possesses sufficient temporal regularity.

Then there exists a constant $C>0$, independent of $h$ and $\Delta t$, such that

\[
\max_{0\le n\le N}
\mathbb E
\left[
\|u(t_n)-u_h^n\|^2
\right]
\le
C
\left(
h^2+\Delta t
\right).
\]

\end{theorem}

\begin{proof}[Sketch of proof]

Subtracting the continuous and discrete formulations leads to an equation for the discrete error $\theta^n$.

Using the coercivity of the bilinear form, the stability estimate, and the approximation properties of the Ritz projection, we obtain

\[
\mathbb E
\left[
\|\theta^n\|^2
\right]
\le
C
\left(
h^2+\Delta t
\right).
\]

Combining this estimate with the interpolation error yields the desired result.

\end{proof}

The above theorem shows that the proposed finite element approximation converges strongly in mean square as

\[
h,\Delta t \rightarrow 0.
\]